\documentclass[10pt,a4paper]{amsart}
\usepackage[margin=19mm]{geometry}
\usepackage{amsmath,amssymb,mathtools}
\usepackage[scr=rsfs,cal=euler]{mathalfa}
\usepackage{xfrac}
\usepackage[unicode,hidelinks,bookmarks=false]{hyperref}
\newcommand{\ie}{i.e.\ }
\newcommand{\cf}{cf.\ }
\newcommand{\resp}{respectively\ }
\newcommand{\Subsection}{\subsection}
\providecommand{\nobreakdash}{\nobreak\hbox{-}}
\setlength{\emergencystretch}{2em}
\multlinegap0pt
\usepackage{tikz}
\usepackage{amssymb}
\newtheorem{lem}{Lemma}
\renewcommand{\epsilon}{\varepsilon}
\newcommand{\supp}{{\operatorname{supp\,}}}
\title{Trace estimates and improved pointwise bounds for joint Eigenfunctions}
\author{Xianchao Wu, Xiao Xiao}
\date{Source: arXiv:2604.22197v1 (CC BY 4.0)}
\begin{document}
\maketitle

\noindent\emph{Source and licence.} Trace estimates and improved pointwise bounds for joint Eigenfunctions. Version arXiv:2604.22197v1 (CC BY 4.0), distributed by arXiv under the Creative Commons Attribution 4.0 International licence, http://creativecommons.org/licenses/by/4.0/, as recorded in the arXiv licence metadata for this exact version and shown on the abstract page https://arxiv.org/abs/2604.22197v1. Full licence notice: \texttt{licenses/arxiv-2604.22197-CC-BY-4.0.txt}.\par\smallskip

\noindent\textit{Editorial scope.} Counted unit: the article's lem 2 (source0.tex lines 1048--1053). The article's complete original proof of it is delivered with it (source0.tex lines 1054--1211, one \texttt{proof} environment of 158 lines). No mathematics is written, repaired or re-derived: the statement and the proof are the article's own lines, in the article's own notation, and no retained line is substituted or re-flowed.  Retained uncounted as the source-local context the proof needs: lines 727--738.  Nothing was omitted from any retained span, so the package carries no omission record.  No retained line contains a citation, so the wrapper carries no bibliography and no bibliography entry.  No retained line contains a figure environment, an image-inclusion command or drawing code, so no image is delivered and the package declares no asset dependency.  Editorial formatting only, in addition: an AMS \texttt{amsart} preamble that restates the article's own declarations and macros used by the retained lines, namely the environments \texttt{lem} on separate counters, together with the standard packages those lines need.  The retained environments print in this document as Lemma 1 in order of appearance, whereas the article prints the same items with the numbering of its own full document.  The complete native source of the article as distributed in the arXiv e-print for this exact version is bundled unchanged as \texttt{sources/199076/source0.tex}.
\begin{lem}\label{lwl}
Assume $A(x, hD)\in \Psi^0(M)$ with symbol $a=\sum_{j=0}^\infty a_j h^j$. Let $\rho\in C_0^\infty(\mathbb{R})$ vanish for $|t|\geq \epsilon_0$, $\epsilon_0>0$ small enough, and satisfy $\rho(0)=1$. Let $\Gamma:=\mathcal{C}_x\cap \supp a$. Then if $(p_1|_{\Gamma}, p_2|_{\Gamma})$ is {\bf rank $1$} at $x$, one has that
\begin{equation}\label{main lemma eq2}
\int_{\mathbb{R}} \hat{\rho}(t) \big(A e^{\frac{i}h t(P_1-E_1)}  \rho\big(h^{-1}[P_2(h)-E_2]\big)  A^*\big)(t,x,x) dt =O(h^{-n+2}).
\end{equation}

Otherwise, one has the following
\begin{equation}\label{main lemma eq}
\int_{\mathbb{R}} \hat{\rho}(t) \big(A e^{\frac{i}h t(P_1-E_1)}  \rho\big(h^{-1}[P_2(h)-E_2]\big)  A^*\big)(t,x,x) dt=O(1) I_h(x)^2,
\end{equation}
if $(p_1, p_2)$ is Morse type.
\end{lem}

\begin{lem}\label{lwl 2}
Assume $A(x, hD)\in \Psi^0(M)$ with symbol $a=\sum_{j=0}^\infty a_j h^j$. Let $\rho\in C_0^\infty(\mathbb{R})$ vanish for $|t|\geq \epsilon_0$, $\epsilon_0>0$ small enough, and satisfy $\rho(0)=1$. Let $\Gamma:=\mathcal{C}_x\cap \supp a$. Then if $(p_1|_\Gamma, p_2|_\Gamma, p_3|_\Gamma)$ is {\bf rank $2$} at $x$, one has that
\begin{equation}\label{main lemma 2 eq}
\int_{\mathbb{R}} \hat{\rho}(t) \big(A e^{\frac{i}h t(P_1-E_1)}  \rho\big(h^{-1}[P_2(h)-E_2]\big)\rho\big(h^{-1}[P_3(h)-E_3]\big)  A^*\big)(t,x,x) dt =O(h^{-n+3}).
\end{equation}
\end{lem}

\begin{proof}
With repeating the begining of the proof of Lemma \ref{lwl}, one need to show that the right hand of \eqref{main lemma 2 eq} can bound
\begin{align}
I(x,x;h):=h^{- 6n}\int &\hat{\rho}(s_2)\hat{\rho}(s_1)\hat{\rho}(t) \exp[\frac{i}h \left<x-y_1, \xi\right>] \exp[\frac{i}h \Phi(s_1, s_2,t,y_1, w_1, w_2, y_2, \xi_1, \xi_2, \xi_3)]\cdot  \nonumber \\
& \exp[\frac{i}h \left<y_2-x, \eta\right>] c_0(s_1, s_2,t,y_1,w_1,w_2,y_2,\xi_1,\xi_2,\xi_3;h) a(x, \xi) \cdot \nonumber \\
 &\qquad\qquad \overline{a(x, \eta)} dy_1 d\xi dw_1 d\xi_1 dw_2 d\xi_2 d\xi_3 dy_2 d\eta dt ds_1 ds_2,
\end{align}
here $c_0\sim \sum_{j=0}^\infty c_{0,j} h^j$ with $c_{0,j}\in C^\infty$ and $c$ is supported near the energy shell $\{p_1=E_1\}\cup\{p_2=E_2\}\cup\{p_3=E_3\}$, 
\begin{align*}
\Phi(\cdot)=&\left<y_1-w_1,\xi_1 \right>+\left<w_1-w_2, \xi_2\right>+\left<w_2-y_2, \xi_3\right> +t\big(p_1(y_1,\xi_1)-E_1\big)+s_1\big( p_2(w_1,\xi_2)-E_2 \big)\nonumber \\
&+s_2\big( p_3(w_2,\xi_3)-E_3 \big)+O_{y_1,\xi_1}(t^2)+O_{w_1,\xi_2}(s_1^2)+O_{w_2,\xi_3}(s_2^2).
\end{align*}

Now apply stationary phase in $(y_1, \xi)$-variables. The critical point equations for $(y_1, \xi)$ are
\begin{align}
\xi&=\xi_1+ t\partial_{y_1} p_1(y_1,\xi_1)+ O(t^2), \\
y_1&=x.
\end{align}

With applying Taylor expansion at the critical point $(y_{1c}, \xi_{c})$ one has that
\begin{align}
&I (x,x; h):=h^{-4n}\int \hat{\rho}(s_2)\hat{\rho}(s_1)\hat{\rho}(t)  \exp\Big[\frac{i}h\Big(\left<x-w_1, \xi_1\right>+\left<w_1-w_2,\xi_2\right>+\left<w_2-y_2,\xi_3\right> \nonumber \\
&\qquad+\left<y_2-x, \eta\right> +t(p_1(x,\xi_1)-E_1\big)+ s_1\big(p_2(w_1,\xi_2)-E_2\big) +s_2\big( p_3(w_2,\xi_3)-E_3 \big)+O_{x,\xi_1}(t^2) \nonumber\\
&\qquad\qquad + O_{w_1,\xi_2}(s_1^2)+O_{w_2,\xi_3}(s_2^2) \Big)\Big]\cdot c_1(s_1, s_2,t,x, w_1,w_2,y_2,\xi_1,\xi_2,\xi_3;h)[a(x, \xi_1)+o(1)]  \nonumber\\
&\qquad\qquad\qquad \cdot \overline{a(x, \eta)} dw_1 d\xi_1 dw_2 d\xi_2 d\xi_3 dy_2 d\eta dt ds_1 ds_2, 
\end{align}
here $c_1\sim \sum_{j=0}^\infty c_{1,j} h^j$ and $c_{1,j}\in C^\infty$.

Apply stationary phase in $(y_2, \eta)$-variables. The critical point equations for $(y_2, \eta)$ are
\begin{align}
\eta&=\xi_3, \\
y_2&=x.
\end{align}
With applying stationary phase at the critical point $(y_{2c}, \eta_{c})$ and Taylor expansion one has that
\begin{align*}
I (x,x; h)=&h^{-3n}\int \hat{\rho}(s_2)\hat{\rho}(s_1)\hat{\rho}(t)  \exp\Big[\frac{i}h\Big(\left<x-w_1, \xi_1\right>+\left<w_1-w_2,\xi_2\right>+\left<w_2-x,\xi_3\right> \nonumber \\
&\quad +t(p_1(x,\xi_1)-E_1\big)+ s_1\big(p_2(w_1,\xi_2)-E_2\big) +s_2\big( p_3(w_2,\xi_3)-E_3 \big)+O_{x,\xi_1}(t^2) \nonumber\\
&\quad\quad + O_{w_1,\xi_2}(s_1^2)+O_{w_2,\xi_3}(s_2^2) \Big)\Big]\cdot c_2(s_1, s_2,t,x, w_1,w_2,\xi_1,\xi_2,\xi_3;h) [a(x, \xi_1)+o(1)]  \nonumber\\
&\quad\qquad\quad \cdot \overline{a(x, \xi_3)} dw_1 d\xi_1 dw_2 d\xi_2 d\xi_3 dt ds_1 ds_2, 
\end{align*}
here $c_2\sim \sum_{j=0}^\infty c_{2,j} h^j$ and $c_{2,j}\in C^\infty$.

Apply stationary phase in $(w_1, \xi_2)$-variables. The critical point equations for $(w_1,\xi_2)$ are
\begin{align}
\xi_2&=\xi_1-s_1\partial_{w_1} p_2(w_1,\xi_2)+ O_{w_1,\xi_2}(s_1^2), \\
%\xi_3&=\xi_2-s_2\partial_{w_2} p_3(w_2,\xi_3)+ O_{w_2,\xi_3}(s_2^2),\\
w_1&=w_2-s_1\partial_{\xi_2} p_2(w_1,\xi_2)+ O_{w_1, \xi_2}(s_1^2).
%w_2&=x-s_2\partial_{\xi_3} p_3(w_2,\xi_3)+ O_{w_2,\xi_3}(s_2^2).
\end{align}

Hence with applying Taylor expansion one can get that
\begin{align*}
I (x,x; h)=&h^{-2n}\int \hat{\rho}(s_2)\hat{\rho}(s_1)\hat{\rho}(t)  \exp\Big[\frac{i}h\Big(\left<x-w_2, \xi_1-\xi_3\right>+t(p_1(x,\xi_1)-E_1\big)+ \nonumber \\
&  s_1\big(p_2(w_2,\xi_1)-E_2\big)+s_2\big( p_3(w_2,\xi_3)-E_3 \big)+O_{x,\xi_1}(t^2)+ O_{w_2,\xi_1}(s_1^2)+O_{w_2,\xi_3}(s_2^2) \Big)\Big] \nonumber\\
&\qquad \cdot c_3(s_1,s_2,t,x,w_2,\xi_1, \xi_3;h)[a(x, \xi_1)+o(1)] \overline{a(x, \xi_3)} d\xi_1 dw_2 d\xi_3 dt ds_1 ds_2,
\end{align*}
here $c_3\sim \sum_{j=0}^\infty c_{3,j} h^j$ and $c_{3,j}\in C^\infty$.

Next apply stationary phase in $(w_2, \xi_3)$-variables. The critical point equations for $(w_2,\xi_3)$ are
\begin{align}
\xi_3&=\xi_1-s_1\partial_{w_2} p_2(w_2,\xi_1)-s_2\partial_{w_2} p_3(w_2,\xi_3)+O_{w_2,\xi_1}(s_1^2)+ O_{w_2,\xi_3}(s_2^2),\\
w_2&=x-s_2\partial_{\xi_3} p_3(w_2,\xi_3)+ O_{w_2,\xi_3}(s_2^2).
\end{align}

So applying Taylor expansion one can get that
\begin{align}
I (x,x; h)=&h^{-n}\int \hat{\rho}(s_2)\hat{\rho}(s_1)\hat{\rho}(t)  \exp\Big[\frac{i}h\Big(t(p_1(x,\xi_1)-E_1\big)+ s_1\big(p_2(x,\xi_1)-E_2\big)+ \nonumber \\
&\qquad  s_2\big( p_3(x,\xi_1)-E_3 \big)+O_{x,\xi_1}(t^2)+O_{x,\xi_1}(s_1 s_2)+ O_{x,\xi_1}(s_1^2)+O_{x,\xi_1}(s_2^2) \Big)\Big] \nonumber\\
&\qquad\qquad \cdot c_4(s_1,s_2,t,x,\xi_1;h) [a(x, \xi_1)+o(1)] [\overline{a(x, \xi_1)}+o(1)] d\xi_1 dt ds_1 ds_2,
\end{align}
here $c_4\sim \sum_{j=0}^\infty c_{4,j} h^j$ and $c_{4,j}\in C^\infty$.

Now use the geodesic normal coordinate about $x$ and make the change of variables $\xi_1=r\Theta$, where $r=|\xi_1|$ and $\Theta=(\theta_1,\theta_2,\dots,\theta_{n-1})\in S^{n-1}$. One can get
\begin{align}\label{tilde I'}
I (x, x; h)=& h^{-n}\int \hat{\rho}(s_2)\hat{\rho}(s_1)\hat{\rho}(t)  \exp\Big[\frac{i}h\Big( t (r^2-1\big)+s_1\big(p_2(x, r\Theta)-E_2\big)+s_2\big( p_3(x,r\Theta)-E_3 \big) \nonumber\\ 
&+O_{x,\xi_1}(t^2)+O_{x,\xi_1}(s_1 s_2)+ O_{x,\xi_1}(s_1^2)+O_{x,\xi_1}(s_2^2) \Big)\Big]c_4(s_1,s_2,t, x,r\Theta;h)\nonumber \\
&\quad\cdot[a(x, r\Theta)+o(1)] [\overline{a(x, r\Theta)}+o(1)] r^{n-1} dr d\Theta dt ds_1 ds_2.
\end{align}

Again one would like to apply stationary phase in $(r, t)$-variables. The critical point $(r_c,t_c)$ satisfies that
\begin{align}
2r t  +s_1\Theta\cdot \partial_{\xi_1} p_2(x, r\Theta)+s_2\Theta\cdot \partial_{\xi_1} p_3(x, r\Theta)+O_{x,\xi_1}(t^2)+O_{x,\xi_1}(s_1 s_2)\qquad\qquad & \nonumber \\
+ O_{x,\xi_1}(s_1^2)+O_{x,\xi_1}(s_2^2)&=0, \label{2r}\\
(r^2-1)+O_{x,\xi_1}(t)&=0. \label{2E_1}
\end{align}

Taking $\epsilon_0$ small enough in the definition of $\rho$ and combining \eqref{2r} and \eqref{2E_1}, we can get that
\begin{equation}\label{t_c 2}
t_c=O(s_1)+O(s_2).
\end{equation}


Taking derivative about $\theta_j$ in \eqref{2r}, \eqref{2E_1} and combining \eqref{2E_1} one can derive that
\begin{equation}\label{d t_c 2}
\partial_{\theta_j} t_c =O(s_1)+O(s_2).
\end{equation}

Now performing stationary phase at the critical point $\big(r_c, t_c \big)$ in \eqref{tilde I'} and Taylor expansion gives,
\begin{align}
I (x, x; h)=&h^{-n+1}\int \hat{\rho}(s_1)\hat{\rho}(s_2)  \exp\Big[\frac{i}h \Big(s_1\big(p_2(x, \Theta)-E_2\big)+s_2\big(p_3(x, \Theta)-E_3\big)+O_{x,\Theta}(s_1t_c) \nonumber\\ 
& +O_{x,\Theta}(s_2t_c)+O_{x,\Theta}(t_c^2)+O_{x,\Theta}(s_1 s_2)+O_{x,\Theta}(s_1^2)+O_{x,\Theta}(s_2^2)\Big) \Big] c_5(s_1, s_2, x, \Theta; h)\nonumber\\
&\cdot [a(x, \Theta)+o(1)] [\overline{a(x, \Theta)}+o(1)] d\Theta ds_1 ds_2,
\end{align}
here $c_5\sim \sum_{j=0}^\infty c_{5,j} h^j$ and $c_{5,j}\in C^\infty$.

Then one can separate the integration
\begin{align*}
\Big| I (x, x; h) \Big|& \leq \Big|h^{-n+1}\int_{|s_1|,|s_2|\leq h} \int \hat{\rho}(s_1)\hat{\rho}(s_2)  \exp[\frac{i}h\Phi(s_1,s_2,z,\Theta) ] \cdot c_6(s_1, s_2, x, \Theta; h)d\Theta ds_1 ds_2 \Big| \nonumber \\
+&\Big|h^{-n+1}\int_{|s_1|\leq h,h<|s_2|<1}\int \cdot\, d\Theta ds_1ds_2\Big|+\Big|h^{-n+1}\int_{h<|s_1|<1,|s_2|\leq h}\int \cdot\, d\Theta ds_1ds_2\Big| \nonumber\\
+& \Big|h^{-n+1}\int_{h<|s_1|,|s_2|<1}\int \hat{\rho}(s_1)\hat{\rho}(s_2)  \exp[\frac{i}h\Phi(s_1,s_2,z,\Theta) ] \cdot c_6(s_1,s_2, x, \Theta; h)d\Theta ds_1ds_2\Big| \nonumber\\
&:=\Big| {I}_1 (x; h)\Big| + \Big| {I}_2 (x; h) \Big|+ \Big| {I}_3 (x; h) \Big|+ \Big| {I}_4 (x; h) \Big|,
\end{align*}
here we set $\Phi(s_1,s_2,x,\Theta)= s_1\big(p_2(x, \Theta)-E_2\big)+s_2\big(p_3(x, \Theta)-E_3\big)+O_{x,\Theta}(s_1t_c) +O_{x,\Theta}(s_2t_c)+O_{x,\Theta}(t_c^2)+O_{x,\Theta}(s_1 s_2)+O_{x,\Theta}(s_1^2)+O_{x,\Theta}(s_2^2)$ and $c_6(s_1,s_2, x, \Theta; h)=c_5(s_1, s_2, x, \Theta; h) [a(x, \Theta)+o(1)] [\overline{a(x, \Theta)}+o(1)]$.

We only focus on $\Big| {I}_1 (x; h) \Big|$ and $\Big| {I}_4 (x; h) \Big|$ since the rest can be dealt with similarly. It's straightforward to get that 
\begin{equation}\label{last1'}
\Big| {I}_1 (z; h)\Big|=O(h^{-n+3}).
\end{equation}


For $\Big| {I}_4 (x; h) \Big|$, %one sets
%\begin{equation*}
%\Phi(s_1,s_2,x,\Theta)=\Phi_1(s_1,s_2,x,\Theta)+\Phi_2(s_1,s_2,x,\Theta),
%\end{equation*}
%where
%\begin{align*}
%\Phi_1(s_1,s_2,x,\Theta)&=s_1\big(p_2(x, \Theta)-E_2\big)+O_{x,\Theta}(s_1t_c) +O_{x,\Theta}(t_c^2)+O_{x,\Theta}(s_1 s_2)+O_{x,\Theta}(s_1^2)\\
%\Phi_2(s_1,s_2,x,\Theta)&=s_2\big(p_3(x, \Theta)-E_3\big) +O_{x,\Theta}(s_2t_c)+O_{x,\Theta}(t_c^2)+O_{x,\Theta}(s_2^2)
%\end{align*}
with the help of \eqref{t_c 2} and \eqref{d t_c 2} the key observation is that for some $1\leq j\neq k\leq n-1$
\begin{align}
\big| \partial_{\theta_j} \Phi(s_1,s_2,x,\Theta) \big|&\geq|s_1| \big(\big|\partial_{\theta_j} p_2(x, \Theta)\big|-O(\epsilon_0)\big)>C_1|s_1|>0, \label{rank2 1}\\
\big| \partial_{\theta_k} \Phi(s_1,s_2,x,\Theta) \big|&\geq|s_2| \big(\big|\partial_{\theta_k} p_3(x, \Theta)\big|-O(\epsilon_0)\big)>C_2|s_2|>0 \label{rank2 2}
\end{align}
due to the {\bf rank $2$ condition}  and the definition of $\rho$.

Also with the help of \eqref{t_c 2} and \eqref{d t_c 2} one has that
\begin{equation}\label{s1s2}
\big|\partial_{\Theta}^\alpha \Phi(s_1, s_2, x,\Theta)\big|=O(s_1)+O(s_2), \quad\text{for a given multi-index } \alpha.
\end{equation}


Hence, one can adapt the argument from the end of Lemma \ref{lwl}: integrate by parts twice with respect to variables $\theta_j$ and $\theta_k$, then apply equation \eqref{s1s2} to conclude that
\begin{align}\label{last2'}
\Big| {I}_4 (x; h) \Big| \leq C h^{-n+5}  \Big|\int_{h<|s_1|,|s_2|<1}\frac{1}{s^2_1 s^2_2}  ds_1 d_2 \Big|=O(h^{-n+3}) .
\end{align}
%\begin{align}\label{last2}
%\Big| {I}_4 (x; h) \Big|&=h^{-n+2} \Big|\int_{h<|s_1|,|s_2|<1} \int \frac{1}{\partial_{\theta_j} \Phi} \big(\partial_{\theta_j} \exp[\frac{i}h  \Phi(s_1,s_2,x,\Theta)]\big)  \cdot c_6(s_1,s_2, x, \Theta; h) d\Theta ds_1 ds_2\Big| \nonumber\\
%&\leq h^{-n+2}\Big|\int_{h<|s_1|,|s_2|<1} \int \frac{1}{\partial_{\theta_j} \Phi} \exp[\frac{i}h  \Phi(s_1,s_2,x,\Theta)]  \cdot \partial_{\theta_j} c_6(s_1,s_2, x, \Theta; h) d\Theta ds_1 ds_2\Big| \nonumber\\
%&+ h^{-n+2} \Big|\int_{h<|s_1|,|s_2|<1} \int \frac{\partial^2_{\theta_j}\Phi }{|\partial_{\theta_j} \Phi|^2}   \exp[\frac{i}h  \Phi(s_1,s_2,x,\Theta)] c_6(s_1, s_2,x, \Theta; h) d\Theta ds_1 ds_2\Big| \nonumber\\
%&=h^{-n+3}\Big|\int_{h<|s_1|,|s_2|<1} \int \frac{1}{\partial_{\theta_j} \Phi} \exp[\frac{i}h  \Phi(s_1,s_2,x,\Theta)]  \cdot \partial_{\theta_j} c_6(s_1,s_2, x, \Theta; h) d\Theta ds_1 ds_2\Big| \nonumber\\
%&+ h^{-n+3} \Big|\int_{h<|s_1|,|s_2|<1} \int \frac{\partial^2_{\theta_j}\Phi }{|\partial_{\theta_j} \Phi|^2}   \exp[\frac{i}h  \Phi(s_1,s_2,x,\Theta)] c_6(s_1, s_2,x, \Theta; h) d\Theta ds_1 ds_2\Big| \nonumber\\
%&\leq C h^{-n+3}  \Big|\int_{h<|s_1|,|s_2|<1}\frac{1}{s_1 s_2}  ds_1 d_2 \Big|=Ch^{-n+3} |\ln h|^2.
%\end{align}



\end{proof}

\end{document}
